\section{A round, step by step}\label{app:round}

\Cref{fig:round-a,fig:round-b,fig:round-c} work one round of the model through on six agents A--F, with the numbers the
equations of \cref{sec:model} produce ($\kappa=0.2$, $T=0.3$, $\varepsilon=1$, $\rho=2.6$, $\delta=10$). Arrows run from
speaker to listener, the direction news travels, with width proportional to attention. A is the seeded agent this
round; E is the listener followed in detail. The companion website shows the same steps as a carousel.

\begin{figure}[p]\centering
\includegraphics[width=0.92\linewidth]{round-example-0}\par\medskip
\includegraphics[width=0.92\linewidth]{round-example-1}\par\medskip
\includegraphics[width=0.92\linewidth]{round-example-2}
\caption{Steps 0--2: the state, where the news appears, and how close everyone is to it.}\label{fig:round-a}
\end{figure}

\begin{figure}[p]\centering
\includegraphics[width=0.92\linewidth]{round-example-3}\par\medskip
\includegraphics[width=0.92\linewidth]{round-example-4}
\caption{Steps 3--4: E scores its incoming links and commits its reward.}\label{fig:round-b}
\end{figure}

\begin{figure}[p]\centering
\includegraphics[width=0.92\linewidth]{round-example-5}\par\medskip
\includegraphics[width=0.92\linewidth]{round-example-6}
\caption{Steps 5--6: E blends the reward into its attention (shown for the whole network after one round), and the hub
that thirty rounds with the same source produce.}\label{fig:round-c}
\end{figure}
